\section{The Time Synchronization Mechanism}
One of the most challenging parts of this thesis has been the development of a time synchronization mechanism for the ST devices, as MetaMotion ones already embed their own closed-source implementation. Three main issues should be overcome to do so: first of all, the limited computational power available on the sensor nodes (in particular, the BlueTile ones), the need to implement this system at the higher application layer of the BLE stack (as lower levels are not accessible by the developers), and the need for a compromise between time accuracy and high data rates. 

The main idea is to leverage on the client, that is a Linux computer running a Python script and a fundamental part of our sensors system. It can be used to provide a reference clock for all the data coming from the sensor nodes. To reduce the number of exchanged packets, an application similar to the receiver-to-receiver one could be used, in which a reference broadcasts messages simultaneously to all devices, and then they can use their reception timestamps to synchronize.  

In particular, the system I implemented consists in the client sending packets to all devices at regular intervals, using the \textit{WRITE\_NO\_RESPONSE} property of a GATT characteristic specifically created for this scope on the sensor nodes. That particular property is the fastest way to send data to a BLE device, as no response is necessary and only one packet is sent. As soon as the remote device receives the packet, the relative callback procedure is automatically called by the BLE stack of the node: in that function, the packet is sent back with its original content (a number) incremented by one, and the internal timestamp. 

The client will then receive that information, and associate the received timestamp to the current Unix epoch time expressed in milliseconds using the \textit{time} library of Python 3.10, which returns a precise value (while previous versions of Python weren't able to do so). Then, the successive samples will be timestamped using the available information: 1) the UNIX time obtained in the resynchronization moment \(t_{UNIX}\), 2) the timestamp of the resynchronization moment \(t_R\) and 3) the device timestamp of the sample under analysis. \(t_S\). The result is calculated with Equation \ref{timestamp-formula}.
\begin{equation} \label{timestamp-formula}
T = t_{UNIX} + (t_S - t_R) 
\end{equation}

The resynchronization procedure can be executed at any desired interval by simply modifying the value in the Python script of the client. In the experiments explained in this thesis, the interval is set to 1s, which was able to give good performances when keeping a data rate of 75~Hz. More details in the Results section. 